\documentclass[11pt,a4paper]{article}

\usepackage[a4paper,margin=1in]{geometry}
\usepackage{fontspec}
\IfFontExistsTF{Times New Roman}{\setmainfont{Times New Roman}}{\setmainfont{TeX Gyre Termes}}
\IfFontExistsTF{Menlo}{\setmonofont{Menlo}}{\setmonofont{Latin Modern Mono}}
\usepackage{amsmath,amssymb}
\usepackage{booktabs}
\usepackage{longtable}
\usepackage{array}
\usepackage{enumitem}
\usepackage{graphicx}
\usepackage{fancyhdr}
\usepackage[round,authoryear]{natbib}
\usepackage[hidelinks]{hyperref}
\usepackage{xurl}
\usepackage{microtype}

\setlength{\parindent}{0pt}
\setlength{\parskip}{0.65em}
\setlist[itemize]{leftmargin=1.4em}
\newcommand{\file}[1]{\path{#1}}
\newcommand{\code}[1]{\path{#1}}
\newcolumntype{L}[1]{>{\raggedright\arraybackslash}p{#1}}

\pagestyle{fancy}
\fancyhf{}
\lhead{\small Stage 3 Surrogate Experiment}
\rhead{\small Price / Premium Comparison}
\cfoot{\thepage}
\renewcommand{\headrulewidth}{0.4pt}

\title{\textbf{Stage 3 Price / Premium Surrogate Report}\\
\large Direct Value Versus Positive Continuation-Premium Prediction}
\author{Codex-assisted surrogate modelling review}
\date{June 19, 2026}

\begin{document}
\maketitle

\begin{abstract}
This report reviews the Stage 3 surrogate experiment on the v1 small-grid American option dataset.
Two fixed multilayer perceptrons were compared: a direct value surrogate and a positive-premium
surrogate that predicts continuation premium and reconstructs value as payoff plus premium. The
experiment uses the existing v1 dataset read-only, keeps Delta, Gamma, and boundary fields as
diagnostics only, does not save model weights, and does not train boundary or Greek heads. The
Stage 3 review decision is \code{READY_FOR_BOUNDARY_DIAGNOSTIC_STAGE}.
\end{abstract}

\tableofcontents
\newpage

\section{Purpose of Stage 3}

Stage 3 asks a narrow modelling question: does predicting nonnegative continuation premium help
preserve the American payoff obstacle better than predicting option value directly? The experiment
is intentionally small and fixed. It is not an architecture search, not a production-risk model,
and not a final research conclusion.

\section{Link to v1 Dataset QA}

The source dataset is \file{results/04_surrogate_dataset/v1_small_grid/dataset_v1_small_grid.npz}.
The v1 QA report accepted 288 regimes and 1,498,464 sampled rows for price-surrogate planning. The
solver source remains the validated baseline American CN/PSOR solver, grounded in the finite
difference American option formulation used throughout the validation ladder
\citep{black_scholes_1973,merton_1973,brennan_schwartz_1977,wilmott_howison_dewynne_1995}.

\begin{longtable}{L{0.26\textwidth} L{0.22\textwidth} L{0.40\textwidth}}
\toprule
Split & Rows & Role \\
\midrule
\endhead
Train & 1,051,006 & Training split; preprocessing fit only here. \\
Validation & 98,857 & Regime-level validation split. \\
Test & 223,729 & Regime-level test split. \\
Stress holdout & 124,872 & Held-out stress combinations. \\
\bottomrule
\end{longtable}

\section{Model Families}

\textbf{Direct value model.} The direct baseline uses the seven normalized v1 features and predicts
\code{value_over_K} with a linear output layer.

\textbf{Positive-premium model.} The premium model uses the same seven features plus
\code{payoff_over_K}, predicts \code{premium_over_K}, and applies a Softplus output so the
predicted premium is nonnegative. The value prediction is reconstructed as
\[
\widehat U/K = \Phi/K + \widehat{(U-\Phi)}/K.
\]
The payoff is analytic contract information computed from spot, strike, and option type. It is not
target leakage: it is known before pricing and is used only for reconstruction and obstacle-aware
structure.

\section{Target Policy}

The only training targets are \code{value_over_K} and \code{premium_over_K}. The fields
\code{boundary_spot_over_K}, \code{delta}, and \code{scaled_gamma} remain diagnostic only.
No boundary head, Delta head, Gamma head, or production Greek target is trained in this stage.

\section{Preprocessing Policy}

Both models use \code{StandardScaler} fit on training rows only. Validation, test, and stress
holdout rows do not influence scaling. The train subset is capped deterministically:
\[
\text{seed}=20260619,\quad \text{cap}=250,000,\quad
\text{selected rows}=250,000.
\]
If the train split exceeds the cap, rows are selected with
\code{np.random.default_rng(20260619).choice(train_indices, replace=False)} and sorted before
batching.

\section{Loss Functions and Training}

Both MLPs use mean squared error and Adam with learning rate \(10^-3\). Training used
10 epochs and batch size 8192. No model weights or \code{.pt} files were saved.

\begin{longtable}{L{0.28\textwidth} L{0.24\textwidth} L{0.24\textwidth}}
\toprule
Model & Final train loss & Final validation loss \\
\midrule
\endhead
Direct value MLP & 2.558e-04 & 2.742e-04 \\
Positive-premium MLP & 4.389e-05 & 2.359e-05 \\
\bottomrule
\end{longtable}

\section{Evaluation Protocol}

Metrics are reported by regime-level split and by diagnostic region. Relative errors use a stable
denominator floor of \(10^-4\). Obstacle and negative-premium rates use a near-zero rate tolerance
of \(10^-4\), with violation amounts counted only above \(10^-10\).

\section{Results by Split}

\begin{longtable}{L{0.24\textwidth} L{0.16\textwidth} L{0.16\textwidth} L{0.16\textwidth} L{0.16\textwidth}}
\toprule
Model & Split & Value RMSE & Value MAE & Premium RMSE \\
\midrule
\endhead
direct\_value\_mlp & validation & 0.016558 & 0.013540 & 0.016558 \\
direct\_value\_mlp & test & 0.033351 & 0.023965 & 0.033351 \\
direct\_value\_mlp & stress\_holdout & 0.028907 & 0.023260 & 0.028907 \\
positive\_premium\_mlp & validation & 0.004857 & 0.003073 & 0.004857 \\
positive\_premium\_mlp & test & 0.012269 & 0.007788 & 0.012269 \\
positive\_premium\_mlp & stress\_holdout & 0.009516 & 0.006607 & 0.009516 \\
\bottomrule
\end{longtable}

\section{Direct Value Versus Positive Premium}

The direct model is the plain baseline. The positive-premium model adds American-option structure by
learning the continuation premium and reconstructing value from the analytic payoff. This comparison
therefore tests whether the obstacle-aware representation reduces payoff-obstacle violations while
keeping price accuracy close to the direct value baseline.

\section{Obstacle and Negative-Premium Diagnostics}

\begin{longtable}{L{0.18\textwidth} L{0.20\textwidth} L{0.20\textwidth} L{0.20\textwidth}}
\toprule
Split & Obstacle rate & Negative premium rate & Review \\
\midrule
\endhead
train & 0 & 0 & PASS \\
validation & 0 & 0 & PASS \\
test & 0 & 0 & PASS \\
stress\_holdout & 0 & 0 & PASS \\
\bottomrule
\end{longtable}

\section{Near-Boundary and Strict-Interior Diagnostics}

Boundary-near rows use the v1 diagnostic mask inherited from Ticket 09/10 logic. Strict-interior rows
exclude maturity, payoff-kink, boundary-near, and nonfinite Greek regions. These regions are not
separate training targets; they are evaluation slices.

\begin{longtable}{L{0.22\textwidth} L{0.16\textwidth} L{0.18\textwidth} L{0.15\textwidth} L{0.16\textwidth}}
\toprule
Model & Split & Region & Rows & Value RMSE \\
\midrule
\endhead
direct\_value\_mlp & test & boundary\_near & 2,219 & 0.013360 \\
direct\_value\_mlp & test & strict\_interior & 193,872 & 0.034586 \\
direct\_value\_mlp & stress\_holdout & boundary\_near & 6,341 & 0.023537 \\
direct\_value\_mlp & stress\_holdout & strict\_interior & 103,124 & 0.029066 \\
positive\_premium\_mlp & test & boundary\_near & 2,219 & 0.008631 \\
positive\_premium\_mlp & test & strict\_interior & 193,872 & 0.011874 \\
positive\_premium\_mlp & stress\_holdout & boundary\_near & 6,341 & 0.002972 \\
positive\_premium\_mlp & stress\_holdout & strict\_interior & 103,124 & 0.008976 \\
\bottomrule
\end{longtable}

\section{Figures}

\begin{figure}[htbp]
\centering
\includegraphics[width=0.78\textwidth]{results/05_surrogate_models/price_premium/figures/loss_curves.png}
\caption{Loss curves.}
\end{figure}

\begin{figure}[htbp]
\centering
\includegraphics[width=0.78\textwidth]{results/05_surrogate_models/price_premium/figures/value_prediction_scatter.png}
\caption{Value prediction scatter.}
\end{figure}

\begin{figure}[htbp]
\centering
\includegraphics[width=0.78\textwidth]{results/05_surrogate_models/price_premium/figures/premium_prediction_scatter.png}
\caption{Premium prediction scatter.}
\end{figure}

\begin{figure}[htbp]
\centering
\includegraphics[width=0.78\textwidth]{results/05_surrogate_models/price_premium/figures/error_by_split.png}
\caption{Error by split.}
\end{figure}

\begin{figure}[htbp]
\centering
\includegraphics[width=0.78\textwidth]{results/05_surrogate_models/price_premium/figures/obstacle_violation_comparison.png}
\caption{Obstacle violation comparison.}
\end{figure}


\section{Limitations}

This experiment uses one fixed MLP configuration and a capped training subset. It does not search
architectures, tune hyperparameters broadly, train a boundary head, train Greek heads, or save model
weights. Delta and Gamma remain diagnostic fields. The stress-holdout results are informative but
not a production-risk guarantee.

\section{What Stage 3 Supports}

Stage 3 supports a first comparison between direct value prediction and an obstacle-aware
positive-premium representation. It also verifies that preprocessing can respect regime-level split
separation and that metrics can be reported by split, boundary-near region, and strict-interior
region.

\section{What Stage 3 Does Not Support}

This stage does not approve production deployment, production Greek labels, broad architecture
selection, boundary-head training, or final paper claims. It does not modify the validated solver and
does not regenerate the v1 dataset.

\section{Recommended Next Stage}

If the review decision is \code{READY_FOR_BOUNDARY_DIAGNOSTIC_STAGE}, the next stage may plan a
boundary-diagnostic modelling experiment. If the decision is \code{REVIEW_REQUIRED_BEFORE_BOUNDARY_STAGE},
the next step is human review of the price/premium errors, obstacle diagnostics, and training setup
before adding any new target heads.

\bibliographystyle{plainnat}
\bibliography{reports/03_solver/references}

\end{document}
